%-------------------------
% Cover Letter
% Author : Jonathan U
% License : MIT
%------------------------

\documentclass[letterpaper,11pt]{article}

\usepackage{latexsym}
\usepackage[empty]{fullpage}
\usepackage{titlesec}
\usepackage{marvosym}
\usepackage[usenames,dvipsnames]{color}
\usepackage{verbatim}
\usepackage{enumitem}
\usepackage[hidelinks]{hyperref}
\usepackage{fancyhdr}
\usepackage[english]{babel}
\usepackage{tabularx}
\usepackage{graphicx}
\input{glyphtounicode}


%----------FONT OPTIONS----------
% sans-serif
% \usepackage[sfdefault]{FiraSans}
% \usepackage[sfdefault]{roboto}
% \usepackage[sfdefault]{noto-sans}
% \usepackage[default]{sourcesanspro}

% serif
% \usepackage{CormorantGaramond}
% \usepackage{charter}

% \usepackage{CharisSIL}
\usepackage{txfonts}
\usepackage[T1]{fontenc}


\pagestyle{fancy}
\fancyhf{} % clear all header and footer fields
\fancyfoot{}
\renewcommand{\headrulewidth}{0pt}
\renewcommand{\footrulewidth}{0pt}

% Adjust margins
\addtolength{\oddsidemargin}{-0.5in}
\addtolength{\evensidemargin}{-0.5in}
\addtolength{\textwidth}{1in}
\addtolength{\topmargin}{-.5in}
\addtolength{\textheight}{1.0in}

\urlstyle{same}

\raggedbottom
\raggedright
\setlength{\tabcolsep}{0in}

% Sections formatting
\titleformat{\section}{
  \vspace{-4pt}\bfseries\scshape\raggedright\Large
}{}{0em}{}[\color{black}\titlerule \vspace{-5pt}]

% Ensure that generate pdf is machine readable/ATS parsable
\pdfgentounicode=1

%-------------------------
% Custom commands
\newcommand{\resumeItem}[1]{
  \item\small{
    {#1 \vspace{-4pt}}
  }
}

\newcommand{\resumeSubheading}[4]{
  \vspace{-2pt}\item
    \begin{tabular*}{0.97\textwidth}[t]{l@{\extracolsep{\fill}}r}
      \vspace{-2pt}
      \textbf{#1} & #2 \\
      \vspace{-1pt}
      \textit{\small#3} & \textit{\small #4} \\
    \end{tabular*}\vspace{-7pt}
}

\newcommand{\resumeSubSubheading}[2]{
    \item
    \begin{tabular*}{0.97\textwidth}{l@{\extracolsep{\fill}}r}
      \textit{\small#1} & \textit{\small #2} \\
    \end{tabular*}\vspace{-7pt}
}

\newcommand{\resumeProjectHeading}[2]{
    \item
    \begin{tabular*}{0.97\textwidth}{l@{\extracolsep{\fill}}r}
      \small#1 & #2 \\
    \end{tabular*}\vspace{-7pt}
}

\newcommand{\resumeSubItem}[1]{\resumeItem{#1}\vspace{-4pt}}

\renewcommand\labelitemii{$\vcenter{\hbox{\tiny$\bullet$}}$}

\newcommand{\resumeSubHeadingListStart}{\begin{itemize}[leftmargin=0.15in, label={}]}
\newcommand{\resumeSubHeadingListEnd}{\end{itemize}}
\newcommand{\resumeItemListStart}{\begin{itemize}}
\newcommand{\resumeItemListEnd}{\end{itemize}\vspace{-5pt}}

%-------------------------------------------
%%%%%%  RESUME STARTS HERE  %%%%%%%%%%%%%%%%%%%%%%%%%%%%


\begin{document}
%----------HEADING----------
% \begin{tabular*}{\textwidth}{l@{\extracolsep{\fill}}r}
%   \textbf{\href{http://sourabhbajaj.com/}{\Large Sourabh Bajaj}} & Email : \href{mailto:sourabh@sourabhbajaj.com}{sourabh@sourabhbajaj.com}\\
%   \href{http://sourabhbajaj.com/}{http://www.sourabhbajaj.com} & Mobile : +1-123-456-7890 \\
% \end{tabular*}

\begin{center}
    \textbf{\huge \uppercase{Full Name}} \textnormal{(pronouns - optional)}\\ \vspace{1pt}
    \small 111 Any Street, Victoria BC V8K 1LM $|$ (250) 123-4567 $|$ lawstudent@uvic.ca 
\end{center}

%------DATE & ADDRESSEE----

\today \par 

\vspace{8pt}

Contact Name, Title\par \vspace{-1pt}
Organization Name\par \vspace{-1pt}
Organization Address\par \vspace{-1pt}
City, Province Postal Code\par 

\vspace{8pt}

%--------RECIPIENT & SUBJECT----
Dear Firstname Lastname, 
\par \vspace{8pt}
\textbf{Re: 2027 Summer Student Position}
\par \vspace{8pt}


%-----BODY TEXT-------
Paragraph 1 – Introduction: The first paragraph simply states the position you are applying for and briefly introduces
yourself, including your law school and current academic level (i.e., “I have recently completed my second year of studies
at the University of Victoria Faculty of Law”). Do NOT start your letter with, “My name is …”. When naming the
position you are applying to, include the full name of the employer. You can then use a shortened version or acronym for
the employer’s name for the rest of the letter. State your connection to the geographic location you are applying to. For
students graduating at a different time that might seem “out-of-sync” from other law students, such as Co-op or JD/JID
students, indicate your anticipated date of graduation. JD/JID students may also wish to include a brief description of the
JID program.

\par \vspace{8pt}

Paragraph 2 – Why Are You Interested In This Employer: The second paragraph addresses why you are applying to
this particular employer and how the position fits with your career goals. Do research on the organization (through the
employer’s website, the Canadian Directory of Legal Employers, and networking) and refer to elements such as their
mission or mandate, practice areas, culture, student programs, mentorship opportunities, pro bono work, or diversity
policies. You may also wish to mention a firm’s size and type. For government, be sure to discuss your interest in public
service. Customize this paragraph for each application. Refer to relevant past experiences briefly as proof of your
interest in the things you claim interest in. If you have interacted with the employer, such as visiting their office, attendingan event, or meeting one of their lawyers for an informational interview or “coffee chat”, mention it here. While this
paragraph may feel similar for many employers, try to make it feel personal to each individual employer. Try to ensure the
interest you are expressing is genuine.

\par \vspace{8pt}

Paragraph 3 – What You Have to Offer: The third paragraph connects your past work experience and skills to what the
employer is looking for. Focus on key skills or competencies that are relevant to the employer’s work, or that are listed in
the job posting, drawing from your work, volunteer, extracurricular, and personal experiences without repeating the
content of your resume. For example, if you are applying to a litigation firm, draw attention to your mooting experience.
A cleverly drafted cover letter will present past experiences, law-related or not, in a way that show these transferrable
skills. Don’t just list your skills without context. Consider the difference in the following statements: “I am a team player
with efficient time-management and organizational skills” “Working as a student advisor to over 30 students, I
developed my interpersonal skills and learned how to listen and meet people’s concerns. Furthermore, as a manager at a
local restaurant, my responsibilities included scheduling, assigning tasks, and general employee motivation. From this
experience, I gained valuable human resource experience, developed strong organizational and time management skills,
and fostered my ability to be both a leader and a team player.”

\par \vspace{8pt}

Paragraph 4 – Conclusion: Conclude by thanking the reader for their time and consideration.

\par \vspace{8pt}

%------SIGNATURE-----

Sincerely, \par

\includegraphics[height=1.5cm]{signature.png}

Your Typed name

%-------------------------------------------
\end{document}